\section{الفصل~\ref{sec:gaba}. \nt{GLUT} و\nt{GABA} معًا}

القضايا المنطقية والديناميكية في الفصل مستنتجة فيه نفسه من التحليل الخطي وقانون أوم؛ وتُظهر التجربة أن الدارات التي تقوم عليها، أي الحظر مع التأخير، والتثبيط الأمامي التالي للإثارة، والتثبيت بالتثبيط، وإيقاع حلقة \foreignlanguage{english}{\nt{GLUT}--\nt{GABA}}، تعمل فعلًا في الدماغ.

{\footnotesize
\setlength{\tabcolsep}{3pt}\renewcommand{\arraystretch}{1.15}
\begin{longtable}{>{\raggedright}p{0.5cm}>{\raggedright}p{4.0cm}>{\raggedright}p{5.6cm}>{\raggedright}p{2.3cm}>{\raggedright\arraybackslash}p{1.7cm}}
\toprule
م & القضية & التجربة وما قيس & المصدر & الحالة \\
\midrule\endfirsthead
\toprule
م & القضية & التجربة وما قيس & المصدر & الحالة \\
\midrule\endhead
1 & عصبونات \nt{GABA} في القشرة من الخُمس إلى الربع & صبغ مناعي للغابا في $10$ مناطق من قشرة القرد: نحو $25\,\%$ من العصبونات في $8$ مناطق، وأقل من $20\,\%$ في القشرة البصرية الأولية. ومراجعة لتنوع العصبونات المثبطة في القشرة & \cite{Hendry1987,Markram2004} & مُقاس \\
2 & ما يفعله التثبيط يحدده إلى حد بعيد موضع الاستقرار: التغصن، أو الجسم، أو موضع ولادة النبضة، أو نهاية سلك عصبون آخر & مراجعة: أصناف عصبونات \nt{GABA} في القشرة تختلف بهدفها على المستقبِل. تسجيل متزامن من جسم عصبون هرمي وتغصنه: التثبيط الأمامي أقوى بكثير على الجسم منه على التغصنات. والتثبيط على نهاية سلك عصبون آخر في الحبل الشوكي يقلل قذف الناقل العصبي لخط دخل واحد (مراجعة للقياسات) & \cite{Markram2004,Pouille2001,Rudomin1999} & مُقاس \\
3 & قلب العلامة عبر وسيط يكلف تأخيرًا واحدًا، نحو جزء من الألف من الثانية؛ والتثبيط التالي للإثارة يضيّق نافذة التزامن & العصبونات الهرمية في حُصين الجرذ في الشرائح: التثبيط عبر وسيط واحد يصل عقب الإثارة المباشرة بتأخير قصير، والنافذة التي يُجمع فيها مدخلان مثيران حتى العتبة أقل من $2\ms$. وعلى التغصنات، حيث هذا التثبيط أضعف، النافذة أوسع & \cite{Pouille2001} & مُقاس \\
4 & الحظر $a\wedge\bar b$~\eqref{eq:veto} مع \foreignlanguage{english}{OR} وواحد ثابت تامّ وظيفيًا (القضية~\ref{prop:gabacomplete}) & البرهان في الفصل. نتيجة كلاسيكية: شبكة من العناصر العتبية ذات المداخل المثبطة تحقق أي تعبير منطقي. قضية عن النموذج، ولا تجربة & \cite{McCullochPitts1943} & نظرية \\
5 & الحظر مع التأخير يعطي كاشفًا لاتجاه الحركة بسرعة مثلى $v^*=\Delta x/d$~\eqref{eq:vstar} & شبكية الأرنب: تنشأ الانتقائية للاتجاه بتثبيط يطفئ الإثارة عند الحركة في الاتجاه ”الصفري“. وأظهر التسجيل المنفصل للمدخلين المثير والمثبط للخلايا الانتقائية أن الخلايا عديمة المحاور النجمية (\nt{GABA}) تثبطها تثبيطًا غير متناظر، بالزوائد المتجهة إلى الجهة ”الصفرية“ فقط. والصيغة $v^*$ استنتاج الفصل & \cite{Barlow1965,Fried2002} & مُقاس \\
6 & التثبيط التحويلي يقسم الجهد ولا يطرح منه~\eqref{eq:shunt} (الشكل~\ref{fig:shunt}) & قانون أوم لمقسّم من ثلاث موصليات مع جهد اتزان للكلور قرب الراحة؛ لعصبون لا يطلق نبضات & استنتاج في الفصل & نظرية \\
7 & يؤثر التحويل في تردد النبضات بالطرح، أما القسمة فتنتج عن التثبيط مع الضجيج الخلفي المتوازن & النموذج: في العصبون المنطلق يكاد التيار عبر التحويل يكون ثابتًا، والأثر في التردد طرحي. التجربة: أُدخلت في العصبونات الهرمية للقشرة الحسية الجسدية لدى الجرذ (شرائح) موصليات خلفية مثيرة ومثبطة عبر المشبك الديناميكي؛ والنمو المتوازن المتزامن للاثنتين يقسم ميل الاستجابة دون إزاحة ملحوظة للتردد. مراجعة: القسمة على الفعالية الكلية توجد في مناطق كثيرة من الدماغ & \cite{HoltKoch1997,Chance2002,Carandini2012} & مُقاس \\
8 & عند $\beta>1$ يختار التثبيط المتبادل فائزًا واحدًا (القضية~\ref{prop:wta}، \eqref{eq:wta}) & القيم الذاتية $-(1\pm\beta)/\tau$ استنتاج في الفصل. والترددان $0.73$ و$0.53$ عند $\beta=0.5$ نقطة ثابتة $r_{1,2}=(I_{1,2}-\beta I_{2,1})/(1-\beta^2)$؛ ولا برنامج في \texttt{models/}. ولا يورد الفصل تجربة مباشرة & استنتاج في الفصل & نظرية \\
9 & تصفير الذاكرة بإشارة عبر \nt{GABA}؛ ومجموعتان بتثبيط متبادل تشكلان ماسكًا & ينتج من الشكل~\ref{fig:bistable} ومن القضية~\ref{prop:wta}. لا تجربة & استنتاج في الفصل & نظرية \\
10 & اتزان حلقة \foreignlanguage{english}{\nt{GLUT}--\nt{GABA}} مستقر إذا كان التثبيط قويًا بما يكفي وسريعًا بما يكفي (القضية~\ref{prop:ei}) & نموذج الجمهرتين~\eqref{eq:ei}؛ والشرطان (i) و(ii) هما محدد مصفوفته وأثرها، مستنتجان في الفصل & \cite{WilsonCowan1972} & نظرية \\
11 & عند $w_{EE}=1.5$ و$w_{EI}w_{IE}=2$ يحفظ التثبيط السريع ($\tau_I=2\ms$) المستوى $E^*=0.67h$، والبطيء ($30\ms$) يؤرجح التذبذبات (الشكل~\ref{fig:ei}) & حل عددي لـ~\eqref{eq:ei}، البرنامج \texttt{figures/make\_ei.py} & حساب & نموذج \\
12 & تعمل القشرة في نمط التثبيت بالتثبيط؛ وبصمته: ادفع عصبونات \nt{GABA} فيهبط ترددها~\eqref{eq:isn} & تنبأت النظرية بالاستجابة المتناقضة. التجربة: تسجيلات داخل الخلية في القشرة البصرية الأولية للقط: عند كبت الاستجابة بمنبه محيط لا تزداد الموصلية المثبطة بل تهبط مع المثيرة. والإثارة البصرية الوراثية لعصبونات \foreignlanguage{english}{PV} لدى الفأر في القشرة البصرية والحسية الجسدية والحركية تخفض تردد العصبونات المثبطة، حتى من دون منبه وتحت تخدير خفيف. والمعامل $-1/3$ بأرقام الشكل~\ref{fig:ei} استنتاج الفصل & \cite{Tsodyks1997isn,Ozeki2009,Sanzeni2020} & مُقاس \\
13 & الحلقة ”\nt{GLUT} يثير \nt{GABA}، و\nt{GABA} يثبط \nt{GLUT}“ تولّد إيقاعًا سريعًا دوره $12\text{--}50\ms$ & التحفيز البصري الوراثي لعصبونات \nt{GABA} السريعة في القشرة الحسية الجسدية للفأر بترددات $8\text{--}200\,\text{Hz}$ يقوّي انتقائيًا إيقاع غاما ($20\text{--}80\,\text{Hz}$، الدور $12\text{--}50\ms$)؛ وتحفيز عصبونات \nt{GLUT} لا يقوّي إلا الإيقاعات الأبطأ & \cite{Cardin2009,Buzsaki2012} & مُقاس \\
\bottomrule
\end{longtable}}
